% Section 1: Introduction
\section{Introduction}

Across 223 large language models evaluated on AlpacaEval 2.0, model capability predicts
length-controlled preference with a partial Spearman correlation of 0.985 --- \textit{after
fully accounting for response verbosity}. Two independent evaluation systems, one based on
human pairwise comparison and one on AI-debiased annotation, agree on model capability ordering
to a degree the field has not previously quantified. More strikingly, this agreement is stronger
once verbosity is removed than before: the bivariate correlation reported by
\citet{Dubois2024LengthControlled} is 0.94; our partial correlation is 0.985. Controlling for
the confound actually reveals the signal more clearly.

This finding bears on a question at the center of LLM evaluation: does length-controlled win
rate (\texttt{LC\_winrate}) preserve the capability ordering that human-annotated win rate
(\texttt{win\_rate}) reflects? The LC correction in AlpacaEval 2.0 is a GLM-based debiasing
procedure that regresses out response length from pairwise comparison outcomes
\citep{Dubois2024LengthControlled}. Its stated goal is to remove verbosity inflation --- the
tendency of AI judges to favor longer, more elaborate responses --- while leaving genuine
quality differences intact. Whether it succeeds at preserving capability ordering at population
scale has not been tested with confound-controlled methods.

The surface problem is well-known: verbosity biases AI-based evaluation.
\citet{Zheng2023Judging} and \citet{Shi2024Judging} document that AI judges systematically
prefer longer responses. AlpacaEval 2.0's LC correction is specifically designed to address
this. However, a deeper problem remains unresolved: verbosity and capability co-vary in real
model populations (VIF = 1.764 in our data), meaning that any observed correlation between
\texttt{win\_rate} and \texttt{LC\_winrate} could reflect shared verbosity rather than shared
capability signal. The gap in existing work is precisely this: no study has applied partial
correlation or standardized regression to disentangle the independent contributions of
capability and verbosity to LC preference across a large, diverse model population.

Our key insight is that this is a partial correlation question, not a bivariate one. Once we
control for \texttt{avg\_length}, the capability-LC alignment does not weaken --- it
strengthens. This is because verbosity was acting as a mild suppressor: capable models tend to
write longer responses, which partially masks the true capability-LC signal in bivariate
analysis. After residualization, the underlying capability signal dominates the LC prediction
by a factor of ${\sim}4.9$ (standardized $\beta$ ratio: 21.34 vs $-4.37$), and the negative
$\beta$ for verbosity confirms the LC correction is functional.

We make the following contributions:

\begin{enumerate}
  \item \textbf{First population-scale partial correlation quantification}: We establish
    $\rho(\texttt{win\_rate}, \texttt{LC\_winrate} \mid \texttt{avg\_length}) = 0.985$
    ($p = 1.69 \times 10^{-170}$ one-tailed, pre-registered directional test, $N=223$),
    with bootstrap 95\% CI [0.976, 0.988].

  \item \textbf{Capability-verbosity dominance via standardized regression}: Capability
    contributes ${\sim}4.9\times$ more to LC preference than verbosity
    ($|\beta_{\mathrm{win}}| = 21.34$ vs $|\beta_{\mathrm{len}}| = 4.37$;
    $R^2 = 0.963$ vs verbosity-only $R^2 = 0.256$ --- 70 percentage points more explained
    variance).

  \item \textbf{FWL theorem verification}: Two independent estimators converge within 1.1pp
    ($\rho_{\mathrm{resid}} = 0.974$, $\delta = 0.011$), ruling out methodological artifact.

  \item \textbf{Quartile-level monotonicity with large effect}: KW $\varepsilon^2 = 0.883$
    ($H = 196.32$, $p = 2.63 \times 10^{-42}$), with fully monotonic LC ordering
    (Q1 = 7.14 $<$ Q2 = 14.69 $<$ Q3 = 26.41 $<$ Q4 = 51.62) and Dunn Q1 vs Q4
    $p = 1.04 \times 10^{-38}$.

  \item \textbf{Methodological lesson on dependent variable choice}: KW on $\Delta$ yields
    $\varepsilon^2 = 0.088$ vs $\varepsilon^2 = 0.883$ for \texttt{LC\_winrate} --- a
    $10\times$ attenuation from DV choice alone, with concrete recommendations for future
    studies.
\end{enumerate}

We organize the paper as follows. Section~\ref{sec:related} reviews prior work.
Section~\ref{sec:methodology} describes methodology. Section~\ref{sec:experiments} presents
experimental design. Section~\ref{sec:results} reports results. Section~\ref{sec:discussion}
discusses limitations. Section~\ref{sec:conclusion} concludes.
